% !TeX program = xelatex
% !TeX encoding = UTF-8
% =====================================================================
% MathorCup（妈妈杯）论文骨架模板 —— 与国赛模板（latex模板/）同一套版式
%
% 版式来源：仓库 latex模板/JXUSTmodeling.cls（字体、字号、行距、章节编号、
%           图表标题、附录代码样式全部沿用），仅按 MathorCup 官方规范补结构。
%
% 官方规范（2026 年第十六届，组委会 2026 年 4 月修订）对应关系：
%   ① 第 1 页：论文题目 + 摘要 + 关键词（本模板 abstract 环境）
%   ② 第 2 页：目录（本模板 \mulu 命令，自动另起一页）
%   ③ 目录之后为正文，页码从 1 起、位于每页页脚中部（\mulu 已自动切换）
%   ④ 全文无页眉、无学校与个人信息（类文件已设 \pagestyle{empty} 前置）
%   ⑤ 正文控制在 30 页左右，附录页数不限
%   ⑥ 参考文献按正文引用顺序编号，正文用方括号标注（类文件 natbib 上标方括号）
%   ⑦ 附录需给出软件名称、运行命令与全部源程序（附录 A、B）
%   ⑧ 使用了 AI 工具须在「参考文献之后」声明（见文末 AI工具使用声明 一节）
%
% 使用说明：
%   1. 编译必须用 XeLaTeX，连跑两遍（或运行 build.bat / build.sh）
%   2. 图片统一放入 figures/ 目录，正文按文件名引用（不带路径与扩展名也可）
%   3. 标注「(待补充：…)」处为按赛题与论文实际情况需要填写的内容
%   4. 正文中标注「示例」的图表仅演示版式，写作时替换为实际内容
%   5. 保留的图表格式清单见同目录《模板说明.md》
% =====================================================================
\documentclass{MathorCupmodeling}
% Please add the following required packages to your document preamble:
\usepackage{longtable}
\usepackage{booktabs}
\usepackage{threeparttable}
\usepackage{multirow}
\usepackage{amsmath}
\usepackage{graphicx}
\graphicspath{{figures/}}
\usepackage{pythonhighlight}
%\usepackage{subfigure}
\usepackage{caption,subcaption}
%\usepackage{ctex}  % 类文件基于 ctexart，ctex 已加载，无需重复 \usepackage{ctex}
\usepackage{xcolor}
\usepackage{listings}
\usepackage{colortbl}
\usepackage{float}
\usepackage{multirow}  % 请确保导言区有这行，MathorCupmodeling.cls 通常已自带

% ---------- 附录代码外观：与国赛模板完全一致 ----------
\lstset{
	backgroundcolor=\color{white},      % 底色白色
	basicstyle=\ttfamily\small,
	breaklines=true,                    % 长代码自动换行
	breakindent=0pt,                    % 换行续行不缩进，保持左对齐
	numbers=left,
	numberstyle=\tiny\color{gray},
	numbersep=6pt,
	frame=single,                       % 加一圈细边框，把代码圈起来
	rulecolor=\color{black},            % 边框黑色
	keywordstyle=\color{blue},
	commentstyle=\color{teal!80},       % 注释：蓝绿色
	stringstyle=\color{red!75},
	showstringspaces=false,
	tabsize=4,
	columns=flexible
}

% 恢复 enumerate 默认缩进（覆盖类文件的 leftmargin=0pt）
\setlist[enumerate,1]{leftmargin=*,labelsep=0.5em,itemsep=0pt,parsep=0pt}

% 中文字体：Windows 下用系统宋体（粗体黑体、斜体楷体）；系统没有时由类文件的
% Fandol 开源字体族接管，无需改动。非 Windows 环境也能直接编译。
\IfFontExistsTF{SimSun}{\setCJKmainfont[BoldFont={SimHei}, ItalicFont={KaiTi}]{SimSun}}{}

\begin{document}

% ==================== 首页信息 ====================
% (待补充：队伍编号 MCxxxxxxx；题号 A/B/C/D/E，须与报名系统内选择的一致)
% 这两项取自官方 Word 模板首页；若所在赛区通知不要求，取消下一行注释即可整块隐藏
% \renewcommand{\mukuai}{}
\duiwu{MC0000000}
\tihao{A}

% (待补充：论文题目；关键词 3~5 个，以"；"分隔)
\biaoti{（论文题目）}
\keyword{关键词一；关键词二；关键词三}

% ==================== 摘要（第 1 页，不编页码） ====================
% (待补充：摘要正文，务必控制在一页以内（本版式约 1000 字，含关键词行），否则目录会被挤到第 3 页)
% 写法：背景一句 → 每问"方法+结果"各一句 → 模型特色一句
\begin{abstract}


\end{abstract}

% ==================== 目录（第 2 页，不编页码） ====================
% \mulu 一次性完成：另起一页出目录 → 再另起一页 → 页码从 1 起、页脚中部
% 目录默认只显示两级标题；需要三级标题时在导言区加 \setcounter{tocdepth}{3}
% (注意：目录须控制在一页内，小节过多时可减少 \subsection 或改用 \subsubsection*)
\mulu

% ==================== 一、问题重述 ====================
\section{问题重述}
\subsection{问题背景}
% (待补充：研究背景)

\subsection{问题提出}
% (待补充：逐条简述各问需要解决的问题)

% ==================== 二、问题分析 ====================
\section{问题分析}

\subsection{问题一的分析}
% (待补充：只写思路，不写结果)

\subsection{问题二的分析}
% (待补充)

\subsection{问题三的分析}
% (待补充)

\subsection{问题四的分析}
% (待补充：按实际问数增删)

\subsection{全文总体框架}
% (待补充：全文各问之间的输入输出关系，可配流程图)

% ==================== 三、模型假设 ====================
\section{模型假设}
\begin{enumerate}[label=\arabic*., wide=0pt, labelsep=0.5em]
	\item % (待补充：假设1，需与"模型的缺点"呼应)
	\item % (待补充：假设2)
	\item % (待补充：假设3)
	\item % (待补充：假设4)
\end{enumerate}

% ==================== 四、符号说明 ====================
% 表格格式 T1：跨页长表 longtable（自动续表头，三线加粗样式）
% 写作时随用随补，示例符号请全部替换
\section{符号说明}
\begin{longtable}{ccc}
	\caption{主要符号说明}
	\label{tab:symbol}\\
	\toprule[1.5pt]
	符号 & 含义 & 单位 \\
	\midrule
	\endfirsthead
	\caption[]{主要符号说明(续)}\\
	\toprule[1.5pt]
	符号 & 含义 & 单位 \\
	\midrule
	\endhead
	\bottomrule[1.5pt]
	\endfoot
	$x$ & 决策变量示例 & $-$ \\
	$f(\cdot)$ & 目标函数示例 & $-$ \\
	$T$ & 时段总数示例 & 个 \\
	$\Delta t$ & 时段步长示例 & $\mathrm{h}$ \\
	$\alpha$ & 步长因子示例 & $-$ \\
	$\beta$ & 折算系数示例 & $-$ \\
\end{longtable}

% ==================== 五、模型准备 ====================
\section{模型准备}

\subsection{数据侧写}
% (待补充：数据总体描述)

% ---------- 图片格式 F1：单图定宽（[H] 强制当前位置） ----------
% 示例图片 figures/KingL.png 为白底黑字占位图，写作时替换为实际图片
\begin{figure}[H]
	\centering
	\includegraphics[width=0.9\textwidth]{KingL.png}
	\caption{单图版式示例（占位图）}
	\label{fig:single_demo}
\end{figure}
% (待补充：图后紧跟分析文字，句式如"由图X可知……")

% ---------- 图片格式 F2：并排子图（模板预留，未启用） ----------
% 如需并排子图，取消以下注释即可使用，图片仍用 figures/KingL.png 占位：
% \begin{figure}[H]
% 	\centering
% 	\begin{subfigure}[b]{0.48\textwidth}
% 		\centering
% 		\includegraphics[width=\textwidth]{KingL.png}
% 		\caption{子图一占位}
% 		\label{fig:sub_demo_a}
% 	\end{subfigure}
% 	\hfill
% 	\begin{subfigure}[b]{0.48\textwidth}
% 		\centering
% 		\includegraphics[width=\textwidth]{KingL.png}
% 		\caption{子图二占位}
% 		\label{fig:sub_demo_b}
% 	\end{subfigure}
% 	\caption{并排子图版式示例（占位图）}
% 	\label{fig:sub_demo}
% \end{figure}

\subsection{误差特征分析}
% (待补充：按需改名或删除本小节；如做预测误差统计，
%  可参考「符号说明」处的 longtable 格式再建一张误差汇总表)

% ==================== 六~九、各问模型建立与求解 ====================
% 每问统一为「模型建立 / 模型求解 / 求解结果与分析」三段骨架，
% 问题一中保留了 Step1~Step4 的求解分解小节样式，可按需套用

\section{问题一模型建立与求解}
% (待补充：本问总起段——解决什么、用什么模型、输出什么)

\subsubsection*{Step1定义决策变量}
% (待补充：变量定义；多变量定义可参考下面的多行公式版式)

% ---------- 公式版式 E1：多行公式块（编号不参与全文编号） ----------
\begin{equation*}
	\left\{
	\begin{aligned}
		x_1 &,\text{ 变量一含义占位};\\
		x_2 &,\text{ 变量二含义占位};\\
		x_3 &,\text{ 变量三含义占位。}
	\end{aligned}
	\right.
\end{equation*}

\subsubsection*{Step2确立约束条件}
% (待补充：逐条约束，每条约束后写现实含义)

% ---------- 公式版式 E2：单个编号公式 ----------
\begin{equation}
	f(\boldsymbol{x})=\sum_{i=1}^{n} w_i x_i,\quad \boldsymbol{x}\in\mathbb{R}^{n}
	\label{eq:demo_obj}
\end{equation}
式(\ref{eq:demo_obj})中，$w_i$ 为权重占位，$x_i$ 为变量占位（公式参数含义写法示例）。

\subsubsection*{Step3构建目标函数}
% (待补充：目标函数与取值理由)

\subsubsection*{Step4模型整合}
% ---------- 公式版式 E3：多行约束整合块 ----------
\begin{equation*}
	\begin{aligned}
		\min_{\boldsymbol{x}} \quad & f(\boldsymbol{x})=\sum_{i=1}^{n} w_i x_i \\
		\text{s.t.} \quad
		& \begin{cases}
			g_j(\boldsymbol{x}) \le 0, \quad j=1,2,\dots,m \quad &\text{(约束含义占位)}\\
			h_k(\boldsymbol{x}) = 0, \quad k=1,2,\dots,p \quad &\text{(约束含义占位)}\\
			\boldsymbol{x}\in\mathbb{R}^{n}
		\end{cases}
	\end{aligned}
\end{equation*}

\subsection{模型求解}
% (待补充：算法选择理由与求解步骤)

\subsection{求解结果与分析}
% (待补充：结果表/图 + "由表X可知/观察图X可知"分析)

% ---------- 表格格式 T2：全边框表格（|c| 列分隔 + \hline + \multicolumn 合并行） ----------
% 示例数据为占位数据，写作时全部替换
\begin{table}[H]
	\centering
	\caption{全边框表格版式示例（占位数据）}
	\begin{tabular}{|c|c|c|c|c|c|}
		\hline
		时段 & 数值A & 时段 & 数值B & 时段 & 数值C \\
		\hline
		第1段 & 1.20 & 第2段 & 3.40 & 第3段 & 2.80 \\
		\hline
		第4段 & 4.15 & 第5段 & 2.05 & 第6段 & 0.75 \\
		\hline
		\multicolumn{2}{|c|}{合计/数值A} & \multicolumn{4}{c|}{14.35} \\
		\hline
	\end{tabular}
	\label{tab:boxed_demo}
\end{table}

\section{问题二模型建立与求解}
\subsection{模型建立}
% (待补充；衔接句写法示例："问题二建立在问题一的基础之上，要求……")

\subsection{模型求解}
% (待补充)

\subsection{求解结果与分析}
% (待补充)

\section{问题三模型建立与求解}
\subsection{模型建立}
% (待补充)

\subsection{模型求解}
% (待补充)

\subsection{求解结果与分析}
% (待补充)

% ---------- 表格格式 T3：标准三线表（booktabs：\toprule/\midrule/\bottomrule） ----------
% 示例数据为占位数据，写作时全部替换
\begin{table}[H]
	\centering
	\small
	\setlength{\tabcolsep}{6pt}
	\caption{标准三线表版式示例（占位数据）}
	\label{tab:booktabs_demo}
	\begin{tabular}{lrrrr}
		\toprule
		方法 & 指标一 & 指标二 & 指标三 & 指标四 \\
		\midrule
		方法A & 0.812 & 12.6 & 3.45 & 0.27 \\
		方法B & 0.905 & 9.8  & 2.91 & 0.31 \\
		方法C & 0.877 & 10.4 & 3.08 & 0.29 \\
		方法D & 0.931 & 8.7  & 2.64 & 0.35 \\
		\bottomrule
	\end{tabular}
\end{table}

% ---------- 表格格式 T4：三线表 + 合并单元格（\multirow 纵向合并 + \multicolumn 横向合并） ----------
% 示例数据为占位数据，写作时全部替换
\begin{table}[H]
	\centering
	\small
	\setlength{\tabcolsep}{8pt}
	\caption{合并单元格三线表版式示例（占位数据）}
	\label{tab:multirow_demo}
	\begin{tabular}{llccc}
		\toprule
		方案 & 指标 & 场景一 & 场景二 & 场景三 \\
		\midrule
		\multirow{2}{*}{方案A} & 指标一 & 81.5 & 76.2 & 88.0 \\
		& 指标二 & 64.3 & 69.7 & 72.1 \\
		\midrule
		\multirow{2}{*}{方案B} & 指标一 & 85.2 & 79.4 & 83.6 \\
		& 指标二 & 70.8 & 66.5 & 74.9 \\
		\midrule
		\multirow{2}{*}{方案C} & 指标一 & 78.9 & 82.3 & 80.7 \\
		& 指标二 & 61.2 & 71.8 & 68.4 \\
		\midrule
		备注 & \multicolumn{4}{c}{三个方案在三个场景下均为占位数据} \\
		\bottomrule
	\end{tabular}
\end{table}

% ---------- 表格格式 T5：带自动换行宽列的三线表（p{} 列 + \addlinespace 行距） ----------
% 适用：某列文字较长需要限定宽度自动换行；示例数据为占位数据
\begin{table}[H]
	\centering
	\small
	\setlength{\tabcolsep}{4pt}
	\caption{自动换行宽列三线表版式示例（占位数据）}
	\label{tab:pcolumn_demo}
	\begin{tabular}{cp{9.4cm}c}
		\toprule
		编号 & 说明（长文本自动换行列） & 数值 \\
		\midrule
		1 & 此处填写较长的说明文字占位：用于演示 p 列在文字过长时自动换行、行内左对齐的版式效果。 & 12.6 \\
		\addlinespace
		2 & 此处填写较长的说明文字占位：说明列宽度由 p{9.4cm} 控制，可按版面自行调整。 & 8.9 \\
		\addlinespace
		3 & 此处填写较长的说明文字占位：记录行之间的 addlinespace 命令用于在三线表中适度增加行间距。 & 10.3 \\
		\bottomrule
	\end{tabular}
\end{table}

\section{问题四模型建立与求解}
% (待补充：按实际问数增删本章)
\subsection{模型建立}
% (待补充)
\subsection{模型求解}
% (待补充)
\subsection{求解结果与分析}
% (待补充)

% ==================== 十、模型检验与灵敏度分析 ====================
\section{模型检验与灵敏度分析}
% (待补充：优化/物理类题写灵敏度分析；统计/机器学习类题写稳健性检验)

% ==================== 十一、模型评价与推广 ====================
\section{模型评价与推广}
\subsection{模型的优点}
\begin{enumerate}[label=\arabic*., wide=0pt, labelsep=0.5em]
	\item % (待补充：每条扣具体机制，禁止空泛)
	\item % (待补充)
	\item % (待补充)
\end{enumerate}

\subsection{模型的缺点}
\begin{enumerate}[label=\arabic*., wide=0pt, labelsep=0.5em]
	\item % (待补充：每条对应一个简化假设，与"模型假设"呼应)
	\item % (待补充)
	\item % (待补充)
\end{enumerate}

\subsection{模型的改进}
\begin{enumerate}[label=\arabic*., wide=0pt, labelsep=0.5em]
	\item % (待补充)
	\item % (待补充)
	\item % (待补充)
\end{enumerate}

\subsection{模型的推广}
\begin{enumerate}[label=\arabic*., wide=0pt, labelsep=0.5em]
	\item % (待补充)
	\item % (待补充)
	\item % (待补充)
\end{enumerate}

% ==================== 参考文献 ====================
% 格式按官方规范：按正文引用顺序编号，正文用方括号标注（类文件已设为上标方括号）
%   书籍：[编号] 作者，书名，出版地：出版社，出版年。
%   期刊：[编号] 作者，论文名，杂志名，卷期号：起止页码，出版年。
%   网络：[编号] 作者，资源标题，网址，访问时间（年月日）。
% 以下为占位示例，写作时按实际引用替换
% 正文引用写法：文献\cite{ref_journal_demo} → 渲染为“文献”+上标[1]（按引用顺序自动编号）
\section*{参考文献}
\renewcommand{\refname}{}
\vspace{-2em}
\begin{thebibliography}{99}
\setlength{\itemsep}{0pt}
	\bibitem[1]{ref_journal_demo}
	作者一，作者二，作者三，等. 文献题名占位，期刊名称占位，45(3)：12-20，2024。
	\bibitem[2]{ref_book_demo}
	作者四. 书名占位，出版地占位：出版社名称占位，2023：55-60。
% 使用了 AI 工具的队伍：取消下一行注释，把工具条目补进参考文献（规范第四条）
%	\bibitem[3]{ref_ai_tool} DeepSeek，DeepSeek-R1-0528，深度求索（DeepSeek），2025-11-25。
\end{thebibliography}

% ==================== AI 工具使用声明 ====================
% 位置规定（MathorCup 与国赛不同，必须注意）：声明放在「参考文献之后」，附录之前，
% 依据《MathorCup 数学应用挑战赛人工智能工具使用规定（试行）》第三条。
% 两种情况二选一：使用了 AI 工具用「情况一」，未使用用「情况二」（可整段替换）。
%
% 情况一：使用了 AI 工具
\section*{AI工具使用声明}
本参赛队在竞赛过程中使用了人工智能工具，具体使用情况如下：
（1）工具名称与版本：DeepSeek，DeepSeek-R1-0528，深度求索（DeepSeek），2025-11-25；
（2）使用环节：建模思路讨论、代码编写与调试、文字润色；
（3）采纳与修改情况：核心模型、算法设计与结果分析由参赛队员独立完成，工具输出内容已逐项核验并修正。
详细的交互记录与采纳说明见支撑材料《AI工具使用详情说明》。
% (待补充：按实际情况增删工具条目；使用了多个工具时逐条列出)
% 注意：使用 AI 工具时，还需在上方参考文献中增列一条工具条目，格式为
%   \bibitem[3]{ref_ai_tool} DeepSeek，DeepSeek-R1-0528，深度求索（DeepSeek），2025-11-25.
%
% 情况二：未使用 AI 工具（使用本情况时，删除上面「情况一」整段，取消下面两行注释）
% \section*{AI工具使用声明}
% 本参赛队未使用任何AI工具。

% ==================== 附录 ====================
\newpage
% 规范：附录须给出实际使用的软件名称、运行命令与全部源程序；附录页数不限
\section*{附录A软件与文件明细}

\subsection*{A.1 软件与运行环境}
% 表格格式 T1 变体：软件清单同样使用跨页 longtable；示例行请按实际替换
\begin{longtable}{ccc}
	\caption{软件与运行环境}
	\label{tab:software}\\
	\toprule[1.5pt]
	软件名称 & 版本 & 用途 \\
	\midrule
	\endfirsthead
	\caption[]{软件与运行环境(续)}\\
	\toprule[1.5pt]
	软件名称 & 版本 & 用途 \\
	\midrule
	\endhead
	\bottomrule[1.5pt]
	\endfoot
	MATLAB & R2023b & 问题一、三的模型求解 \\
	Python & 3.11 & 数据清洗与绘图（numpy、pandas、matplotlib） \\
	\LaTeX{}（XeLaTeX） & TeX Live 2023 & 论文排版 \\
\end{longtable}

\subsection*{A.2 代码文件明细}
% 表格格式 T1 变体：附录文件清单与「符号说明」同一套 longtable
\begin{longtable}{cc}
	\caption{代码文件列表}
	\label{tab:codefile}\\
	\toprule[1.5pt]
	文件名 & 文件描述 \\
	\midrule
	\endfirsthead
	\caption[]{代码文件列表(续)}\\
	\toprule[1.5pt]
	文件名 & 文件描述 \\
	\midrule
	\endhead
	\bottomrule[1.5pt]
	\endfoot
	code/q1.m & MATLAB 代码示例（附录B 由 lstinputlisting 引入） \\
	code/q1.py & Python 代码示例（附录B 由 lstinputlisting 引入） \\
\end{longtable}

\subsection*{A.3 运行命令}
% 规范：附录应写明实际使用的命令；此处列出可直接复现的运行方式
\begin{lstlisting}[frame=single,numbers=none]
% 问题一：MATLAB 环境
>> cd code
>> run('q1.m')

% 问题二、三：Python 环境（依赖见 code/requirements.txt）
$ cd code
$ python q1.py
\end{lstlisting}

\section*{附录B代码}
% 代码不写在正文里，统一用 \lstinputlisting 从外部文件引入；
% 全部代码文件集中放在同目录的 code/ 文件夹中：
%   MATLAB：\lstinputlisting[language=Matlab]{code/q1.m}
%   Python：\lstinputlisting[language=Python]{code/q1.py}
% 代码外观（黑色细边框、左侧行号、关键字蓝色、注释蓝绿色、字符串红色）
% 由导言区 \lstset 统一控制，无需在正文重复设置；
% 换成你自己的代码时，只需替换 code/ 下对应的文件，或修改下面的路径

\subsection*{B.1 MATLAB 代码（q1.m）}
\lstinputlisting[language=Matlab]{code/q1.m}

\subsection*{B.2 Python 代码（q1.py）}
\lstinputlisting[language=Python]{code/q1.py}

\end{document}
